\section{Future Direction: Aligning Resource-Constrained Mutators}
\label{sec:rl}

While the core $E^3$ framework relies on high-fidelity, large-scale model architectures to act as structural mutators, deploying such systems on consumer hardware remains a challenge. Smaller, parameter-efficient models may struggle with structural syntax constraints, frequently producing code that fails compilation or violates JAX's differentiability boundaries. One possible future direction is a reinforcement learning training pipeline using \textbf{Group Relative Policy Optimization (GRPO)}.

Unlike standard RL pipelines that require a memory-intensive critic network (such as PPO), GRPO computes advantages relative to a group of sampled candidate outputs. This could make it suitable for symbolic and code-generation tasks where sample efficiency and memory limits matter.

\subsection{GRPO Formulation and Rewards}
For a given scientific domain and a current best model state $\mathcal{M}_{old}$, the policy model $\pi_\phi$ (initialized from an 8B base model) samples a group of $G$ candidate equations:
\begin{equation}
    \{\mathcal{M}_1, \mathcal{M}_2, \dots, \mathcal{M}_G\} \sim \pi_\phi(\cdot \mid \mathcal{M}_{old})
\end{equation}

Each candidate $\mathcal{M}_i$ is compiled by JAX and evaluated by the ABC-SMC engine to receive a scalar reward $R_i$:
\begin{equation}
    R_i = R_{syntax}(\mathcal{M}_i) + R_{performance}(\mathcal{M}_i) + R_{parsimony}(\mathcal{M}_i)
\end{equation}
where:
\begin{itemize}
    \item $R_{syntax}(\mathcal{M}_i) = +1$ if the Python code successfully compiles and completes the Diffrax simulation without throwing NaNs or memory overflow errors; $-1$ otherwise.
    \item $R_{performance}(\mathcal{M}_i) = \max\left(0, \ln\left(\frac{\text{MSE}_{old}}{\text{MSE}_{new}}\right)\right)$ if the new MSE is lower than the old MSE, providing a continuous, log-scale representation of error reduction.
    \item $R_{parsimony}(\mathcal{M}_i) = -\alpha \cdot \text{length}(\mathcal{M}_i)$ acts as a parsimony penalty intended to discourage high-frequency reward hacking and encourage symbolic sparsity.
\end{itemize}

We compute the relative advantage for each candidate within the sampled group:
\begin{equation}
    A_i = \frac{R_i - \text{mean}(R)}{\text{std}(R) + 10^{-8}}
\end{equation}

The policy parameters $\phi$ are then updated by maximizing the GRPO objective:
\begin{equation}
    L(\phi) = \frac{1}{G} \sum_{i=1}^{G} \min \left( r_i(\phi) A_i, \, \text{clip}(r_i(\phi), 1-\epsilon, 1+\epsilon) A_i \right) - \beta D_{KL}(\pi_\phi \parallel \pi_{ref})
\end{equation}
where $r_i(\phi) = \frac{\pi_\phi(\mathcal{M}_i \mid \mathcal{M}_{old})}{\pi_{ref}(\mathcal{M}_i \mid \mathcal{M}_{old})}$ is the probability ratio, and $D_{KL}$ enforces a Kullback-Leibler divergence penalty to prevent the policy from collapsing away from the reference distribution.

\subsection{Required Evidence}
This GRPO path is not implemented in the current repository. To make it a demonstrated contribution, future work must add training scripts, datasets, model checkpoints or adapters, held-out evaluation data, and syntax/performance metrics that can be reproduced from source.

The primary contribution of the current work remains the LLM-driven structural mutation loop coupled to ABC-SMC evaluation.
